% !TeX root = ../PIFB2.tex
\section{Beyond Mean Field: Correlated Posteriors, the Nested Bound, and Fluctuation Corrections}
\label{sec:beyond_mean_field_gauge_vfe}

Several distinct reductions enter the construction. The factorization $Q=\prod_i q_i$ suppresses state-posterior correlations. A deterministic saddle approximation suppresses fluctuations of whole configurations. An empirical measure is an exact representation of a finite unlabeled population, but replacing its random dynamics by a deterministic large-population law requires an additional limit. Projecting an exact coarse-grained law into a selected functional class is a further approximation. Improving one reduction does not establish the others.

The configuration-space Gibbs lift of Theorem~\ref{thm:belief_configuration_gibbs_lift} places a probability law on whole belief configurations and introduces the distinct configuration temperature $T_{\mathrm{cfg}}$, which is a bookkeeping scale of that ensemble and not a physical temperature. The Laplace term in Eq.~\eqref{eq:rg_laplace} also retains a local Gaussian fluctuation determinant when internal cluster variables are eliminated. Those results do not yet define a correlated posterior over the latent states of different agents, a generating functional for connected configuration correlations, or a self-consistent two-point theory. This appendix states the nested bound that organizes these objects, the correlated state families that preserve evidence semantics, and an outlook on fluctuation corrections, without reclassifying the peer-KL functional as an ordinary state-level evidence bound.

\subsection{A Normalized Nested Variational Model}

Let $X$ denote a generative configuration of prior and transition parameters, supports, graph or source priors, precision fields, and frame or edge variables. It does not contain the variational posterior densities $q_i$ or $s_i$. Let $Y=(Y_1,\ldots,Y_N)$ denote the state-level and model-level latent variables, and let $o$ denote the observations. No slow evolution of $X$ is implied by this probabilistic decomposition. The state model is a fixed conditional law normalized jointly over $(Y,o)$, with likelihood densities relative to one declared observation reference measure
\begin{equation}
p_\theta(dY,o\mid X)
:=p_\theta(dY\mid X)p_\theta(o\mid Y,X),
\label{eq:beyond_mf_state_joint}
\end{equation}
whose factors do not depend on the variational posterior used to approximate them. For a possibly correlated conditional posterior $Q_X(dY)$, define the state variational free energy in nats by
\begin{equation}
\mathcal F_{V,\mathrm{state}}[Q_X;X,o]
:=D_{\mathrm{KL}}\bigl(Q_X\|p_\theta(\cdot\mid X)\bigr)
-\mathbb{E}_{Q_X}\log p_\theta(o\mid Y,X).
\label{eq:beyond_mf_state_vfe}
\end{equation}
It obeys
\begin{equation}
\mathcal F_{V,\mathrm{state}}[Q_X;X,o]
=-\log p_\theta(o\mid X)
+D_{\mathrm{KL}}\bigl(Q_X\|p_\theta(\cdot\mid o,X)\bigr),
\label{eq:beyond_mf_state_elbo_identity}
\end{equation}
so correlated inference changes the approximation family without changing the evidence semantics.

\begin{theorem}[Nested state and configuration VFE]
\label{thm:nested_state_configuration_vfe}
Let $T_{\mathrm{cfg}}>0$, let $P_0(dX)$ be a proper configuration prior, let Eq.~\eqref{eq:beyond_mf_state_joint} be normalized for $P_0$-almost every $X$, and define $p_\theta(o):=\int p_\theta(o\mid X)P_0(dX)\in(0,\infty)$. Let $R\ll P_0$ be a configuration posterior and let $X\mapsto Q_X$ be a measurable probability kernel with $Q_X\ll p_\theta(\cdot\mid X)$ for $R$-almost every $X$. Assume $0<p_\theta(o\mid X)<\infty$ for $R$-almost every $X$. Assume the terms below are finite, or interpret the identity in the extended-real sense without an indeterminate $\infty-\infty$. Then
\begin{align}
\mathcal{J}[R,Q]
&:=T_{\mathrm{cfg}}D_{\mathrm{KL}}(R\|P_0)
+T_{\mathrm{cfg}}\int
\mathcal F_{V,\mathrm{state}}[Q_X;X,o]R(dX)
\nonumber\\
&=-T_{\mathrm{cfg}}\log p_\theta(o)
+T_{\mathrm{cfg}}D_{\mathrm{KL}}
\bigl(R(dX)Q_X(dY)\|P_\theta(dX,dY\mid o)\bigr).
\label{eq:nested_state_configuration_identity}
\end{align}
Consequently, $-\mathcal{J}/T_{\mathrm{cfg}}$ is an evidence lower bound for one normalized hierarchical model, with equality if and only if the joint variational law equals the exact posterior almost everywhere.
\end{theorem}

\begin{proof}
Expand the joint KL divergence by the chain rule, use $P_\theta(dX,dY,o)=P_0(dX)p_\theta(dY,o\mid X)$, and collect the conditional terms into Eq.~\eqref{eq:beyond_mf_state_vfe}.
\end{proof}

The quantity $\mathcal{J}/T_{\mathrm{cfg}}$ is a two-level tower variational free energy with latent $(X,Y)$. Its configuration temperature and the vacuum prior below are specific to this paper. If $P_0$ is itself defined by a vacuum functional,
\begin{equation}
P_0(dX)
:=Z_0^{-1}\exp[-\mathcal{E}_{\mathrm{vac}}(X)/T_{\mathrm{cfg}}]\nu_0(dX),
\label{eq:beyond_mf_configuration_prior}
\end{equation}
assume that $\nu_0$ is a declared sigma-finite reference measure and $0<Z_0<\infty$. For finite terms, the configuration contribution is exactly
\[
T_{\mathrm{cfg}}D_{\mathrm{KL}}(R\|P_0)
=\mathbb E_R\mathcal E_{\mathrm{vac}}
+T_{\mathrm{cfg}}\int\log\frac{dR}{d\nu_0}dR
+T_{\mathrm{cfg}}\log Z_0.
\]
The reference-entropy integral is a nonnegative KL divergence only when $\nu_0$ is a probability law; for Lebesgue reference it is the negative differential entropy and may have either sign. This nests state-level evidence inside the configuration ensemble. The vacuum functional may contain only observation-free structural energies in the declared variables $X$; it cannot contain live $q_i$, $s_i$, or their peer-KL terms without promoting separate measure-valued fields to generative variables and specifying their relation to $Q_X$. The held-out likelihood is excluded from $\mathcal{E}_{\mathrm{vac}}$ and enters once through Eq.~\eqref{eq:beyond_mf_state_vfe}. Environmental variables may be represented as ordinary latent nodes or observed boundary conditions; the construction does not treat a variational belief of another agent as a fixed generative factor.

Both terms of $\mathcal{J}$ carry $T_{\mathrm{cfg}}$, so at the optimal state posterior the optimal configuration posterior is
\[
R^*(dX)\propto
\exp[-\mathcal E_{\mathrm{vac}}(X)/T_{\mathrm{cfg}}]
p_\theta(o\mid X)\nu_0(dX).
\]
A precise concentration statement requires hypotheses on the vacuum energy, reference measure, and likelihood. For example, on a compact configuration space with continuous vacuum energy having a unique minimizer $X_0$, finite reference measure assigning positive mass to each neighborhood of $X_0$, and bounded continuous likelihood positive at $X_0$, the posterior concentrates at $X_0$ as $T_{\mathrm{cfg}}\downarrow0$. To see this, compare the positive energy gap outside any neighborhood with a smaller gap inside a neighborhood on which the likelihood is bounded below. The resulting exterior-to-interior mass ratio decays exponentially. Multiple minima can retain data-dependent relative weights, and a likelihood vanishing on their neighborhoods can exclude vacuum minima entirely. For the separate noncompact Gaussian example $\nu_0(dX)=dX$, $\mathcal E_{\mathrm{vac}}(X)=X^2/2$, and $p_\theta(o\mid X)=\mathcal N(o;X,1)$, completing the square gives
\[
R^*=\mathcal N\left(
\frac{oT_{\mathrm{cfg}}}{1+T_{\mathrm{cfg}}},
\frac{T_{\mathrm{cfg}}}{1+T_{\mathrm{cfg}}}\right).
\]
Here the data move the mean at order $T_{\mathrm{cfg}}$. This scaling does not in general recover minimization of a configuration action in which the data term has an order-one coefficient.

\subsection{Correlated State Posteriors and Region Free Energies}

The first extension beyond mean field preserves the normalized model above while enlarging $Q_X$. For the following marginal-polytope statement, take finite state alphabets and strictly positive finite potentials. Suppose the conditional state prior has a gauge-covariant pairwise factorization on a graph $\mathcal{G}=(V,E)$,
\begin{equation}
p_\theta(Y\mid X)
:=Z_X^{-1}
\prod_{i\in V}\psi_i(Y_i;X)
\prod_{(i,j)\in E}\psi_{ij}(Y_i,Y_j;X),
\label{eq:beyond_mf_pairwise_joint}
\end{equation}
where each edge factor is a fixed model potential and not a function of live variational densities. Introduce node beliefs $b_i$ and edge beliefs $b_{ij}$ satisfying normalization and local marginal consistency. Their Bethe entropy is
\begin{equation}
H_{\mathrm{B}}[b]
:=\sum_{(i,j)\in E}H[b_{ij}]
-\sum_{i\in V}(d_i-1)H[b_i].
\label{eq:beyond_mf_bethe_entropy}
\end{equation}
The corresponding region free energy is the expectation of the fixed node and edge energies $-\log\psi_i$ and $-\log\psi_{ij}$ minus $H_{\mathrm{B}}$. On a tree the local marginal polytope equals the marginal polytope~\cite[Prop.~4.1]{Wainwright2008}, and minimizing this functional over it returns the exact marginals and $-\log Z_X$. On a loopy graph it is the Bethe approximation. Larger regions give the cluster-variation method~\cite{Kikuchi1951} and its region-graph generalizations, a family of more expressive approximations whose accuracy depends on the choice of regions and is not guaranteed to improve~\cite[Sec.~VIII.C]{Yedidia2005}. Generic loopy pseudomarginals need not be the marginals of any global $Q_X$, and their region objective is not a generic upper or lower evidence bound. Bound-preserving constructions require additional structure, such as a globally realizable variational family or an appropriate tree-reweighted relaxation. Ordinary mean field is recovered by the restriction $b_{ij}=b_i b_j$ on every edge when those beliefs arise from a global product law. It is therefore a lower-complexity variational subfamily of the correlated theory, not the definition of the theory.

Gauge covariance acts on the region beliefs as a joint pushforward. If $Y_i\mapsto G_iY_i$ and $Y_j\mapsto G_jY_j$, then $b_{ij}$ is pushed forward by $G_i\times G_j$, while the edge potential is transformed contragrediently so that its scalar energy is unchanged. Entropies must be defined relative to a covariantly transformed or quasi-invariant reference measure, including its Jacobian; a finite invariant reference is not generally available for noncompact $\mathrm{GL}^+(K_q)$. A Regime~I link $\Omega_{ij}=U_iU_j^{-1}$ has trivial loop holonomy and is removed by a passive frame re-expression. A correlated theory with nontrivial holonomy must instead retain the independent links on the interaction graph of Regime~II (Section~\ref{sec:discrete_regime_ii}), with inverse orientation reversal for the reversible group-valued identification sector. Singular selective channels instead form composable maps or Markov kernels, whose cycle compositions are not group holonomies. Integrating only flat vertex-frame identifications cannot generate nontrivial group holonomy.

For the Gaussian sector, a computationally direct correlated family is
\begin{equation}
Q_X(Y)=\mathcal{N}(\bar{Y},C),
\qquad
C^{-1}=\Lambda_Q,
\label{eq:beyond_mf_block_gaussian}
\end{equation}
with a sparse block precision $\Lambda_Q$ supported on the state interaction graph. The off-diagonal blocks $C_{ij}$ carry inter-agent correlations, and if $Y_i=(k_i,m_i)$ contains fast state and slow model variables, the within-node block $C_i^{(km)}$ carries cross-channel correlations that the separate $q_i$ and $s_i$ factors discard. Under a local linear frame change, $C_{ij}\mapsto \rho_i(G_i)C_{ij}\rho_j(G_j)^\top$, where $\rho_i$ is block diagonal if the state and model channels use distinct representations, so the covariance family is closed under the group action. The block-diagonal restriction $C_{ij}=0$ for $i\ne j$ is the Gaussian mean-field limit. At a Gaussian target of precision $\Lambda$, its exact cost at the optimum is the determinant gap $\frac12[\sum_b\log\det\Lambda_{bb}-\log\det\Lambda]\ge0$, which vanishes exactly when all cross-block precisions vanish. Sparse block-Gaussian inference is therefore a tractable extension that exposes a measurable connected two-point state function while retaining tractable linear algebra.

The live alignment term $D_{\mathrm{KL}}(q_i\|\Omega_{ij}q_j)$ does not become an exact instance of Eq.~\eqref{eq:beyond_mf_pairwise_joint} merely because edge beliefs have been introduced: its target still moves with another variational factor, which is the obstruction proved by Theorem~\ref{thm:state_level_elbo_nogo}. It has two coherent placements. It may remain a configuration-level consensus potential, or it may be replaced by a fixed state-transition or pair potential from which the correlated posterior is inferred. Using both for the same interaction would count that edge twice.

\begin{remark}[No-double-counting conditions]
\label{rem:beyond_mf_no_double_counting}
An observation likelihood enters the state free energy once and is not also placed in the vacuum action; a physical interaction is represented either by a fixed edge factor of the state model or by a consensus potential on configurations, not both; a categorical source variable $J_i$ is marginalized exactly or assigned one categorical variational law with probabilities $\beta_{ij}$, and its entropy is counted once; introducing a random simplex-valued parameter is a separate model requiring its own prior and conditional law, not another description of that categorical entropy; and marginal Gaussian entropies and the correlation correction are components of one state-posterior entropy, while attention-row entropy and configuration meta-entropy belong to separate sample spaces. These are instances of two general statements: an energy appended to a joint without its normalizer is not an evidence bound, and the complexity of a correlated recognition law splits exactly into its total correlation and its marginal divergences.
\end{remark}

\subsection{Outlook: Fluctuation Corrections and the Effective Action}

After the state family has been specified, define a state-reduced configuration action
\begin{equation}
\mathcal{A}_o[X]
:=\mathcal{E}_{\mathrm{vac}}[X]
+T_{\mathrm{cfg}}
\inf_{Q_X\in\mathcal{Q}_{\mathrm{corr}}}
\mathcal F_{V,\mathrm{state}}[Q_X;X,o].
\label{eq:beyond_mf_bare_action}
\end{equation}
With the unrestricted posterior family the second term is $-T_{\mathrm{cfg}}\log p_\theta(o\mid X)$; with a normalized family such as a sparse block-Gaussian $Q_X$ it is a variational upper approximation to that conditional surprisal, while a Bethe or Kikuchi pseudomarginal optimization does not inherit this evidence-bound statement. The live multi-agent functional cannot serve as $\mathcal{E}_{\mathrm{vac}}$ in Theorem~\ref{thm:nested_state_configuration_vfe}, because it depends on the variational $q_i,s_i$ and contains the observation term. The peer-KL sector is a gauge-covariant consensus energy and, when its belief-configuration partition function is finite, an exact generalized-Bayesian energy on that separate sample space; inside $\mathcal{A}_o$ it is only a candidate truncation pending an explicit summary map from $Q_X$ to belief fields and quantitative energy and force residuals.

Let $\mathfrak X_{\mathrm{phys}}$ be a regular configuration quotient by transformations explicitly declared to be redundancies; the subscript denotes retained coordinates and does not establish a physical interpretation. Assume a free proper redundancy action on the region used, or work directly on a specified smooth quotient chart. Active transformations remain coordinates of $X$. Choose an injective chart $X=\operatorname{Exp}_{\bar X}\xi$ on a domain $D$ in the tangent space. The following construction describes an ensemble supported in that chart; a global ensemble requires compatible charts and their measure factors, since an exponential map need not be globally injective. For $T=T_{\mathrm{cfg}}>0$, define
\begin{equation}
Z_{\bar X}[J]
:=\int_D
\exp\left[-\frac{\mathcal A_o[\operatorname{Exp}_{\bar X}\xi]-J_A\xi^A}{T}\right]
d\nu_{\mathrm{phys}}(\xi),
\qquad
W_{\bar X}[J]:=T\log Z_{\bar X}[J].
\label{eq:beyond_mf_generating_functional}
\end{equation}
Assume the integral and its first two source derivatives exist on an open neighborhood of the sources considered, with differentiation under the integral justified. Direct differentiation gives
\[
\partial_{J_A}W=\mathbb E_J\xi^A,\qquad
\partial_{J_A}\partial_{J_B}W
=T^{-1}\operatorname{Cov}_J(\xi^A,\xi^B).
\]
Where this covariance is positive definite, the local Legendre transform
$\Gamma_{\bar X}(\varphi)=J_A\varphi^A-W_{\bar X}(J)$ therefore satisfies
\[
\Gamma_{\bar X}^{(2)}
=T\operatorname{Cov}_J(\xi)^{-1}.
\]
At the source-free mean $\varphi_0=\mathbb E_0\xi$ its value is exactly
$\Gamma_{\bar X}(\varphi_0)=-T\log Z_{\bar X}[0]$.
This configuration covariance differs from the state covariance $C_{ij}$ in Eq.~\eqref{eq:beyond_mf_block_gaussian}. The Legendre effective action also differs from the conditional coarse energy $\mathcal E^{\mathrm{exact}}_{\ell+1}$ in Appendix~\ref{sec:meta_agent_rg}, obtained by integrating eliminated variables at a fixed retained value.

A finite-dimensional Laplace statement can now be made precisely. Write
$d\nu_{\mathrm{phys}}=\rho_{\mathrm{phys}}\operatorname{vol}_{g_{\mathrm{phys}}}$.
Assume that $\mathcal A_o$ is $C^6$, has one dominant interior minimum $\bar X$ with positive Hessian $H_{\mathrm{phys}}$, and has a positive energy gap outside every sufficiently small neighborhood of that minimum in a compact integration region. Assume $\rho_{\mathrm{phys}}$ is $C^4$ and positive near $\bar X$. If the action depends on $T$, these regularity, nondegeneracy, and domination hypotheses must hold uniformly. Taylor expansion after $\xi=\sqrt T v$, Gaussian integration, and the cancellation of odd Gaussian moments give
\begin{equation}
\Gamma_{\bar X}(\varphi_0)
=\mathcal A_o(\bar X)
+\frac{T}{2}\log\det\left(
\frac{g_{\mathrm{phys}}^{-1}H_{\mathrm{phys}}}{2\pi T}\right)
-T\log\rho_{\mathrm{phys}}(\bar X)
+O(T^2).
\label{eq:beyond_mf_one_loop}
\end{equation}
The exterior gap makes the omitted exterior integral exponentially small, and the local integral has relative error $O(T)$; multiplication of its logarithm by $T$ gives the stated remainder. In general the mean displacement $\varphi_0$ is $O(T)$ rather than exactly zero, so this equation is evaluated at the source-free mean, not at the classical minimum as an argument of $\Gamma$. An off-stationary effective action requires a source selecting that mean. The factor $2\pi T$ contributes at order $T\log T$ and cannot be suppressed into an $O(T^2)$ remainder. An infinite-dimensional functional determinant additionally needs a regulator, a reference normalization, and a controlled limiting procedure; Eq.~\eqref{eq:beyond_mf_one_loop} alone supplies none of these.

Gauge-fixed integrals require both a transverse Jacobian and the measure induced on the slice. In a finite-dimensional ambient Riemannian space, let a regular gauge condition be $\Phi=0$, let $V_a$ be the redundancy generators, and put
\[
(M_{\mathrm{FP}})_{\alpha a}=d\Phi_\alpha(V_a),
\qquad
J_\Phi=\sqrt{\det(D\Phi g^{-1}D\Phi^\top)}.
\]
The coarea formula gives the local identity
\[
\delta(\Phi)|\det M_{\mathrm{FP}}|\operatorname{vol}_g
=\frac{|\det M_{\mathrm{FP}}|}{J_\Phi}
\operatorname{vol}_{\Phi=0}.
\]
Any ambient density multiplies the right-hand side after restriction to the slice. This is the measure content of the local Faddeev--Popov construction~\cite{FaddeevPopov1967}. Under $\Phi\mapsto A\Phi$ with nonsingular constant $A$, both $|\det M_{\mathrm{FP}}|$ and $J_\Phi$ acquire $|\det A|$, leaving their ratio unchanged. A formula retaining the former determinant alone with the intrinsic slice measure fails this elementary rescaling check. Direct quotient integration and a correctly normalized gauge-fixed representation are alternative descriptions of the same specified measure; neither a second determinant nor an orbit-volume factor can be appended without deriving the relevant disintegration. Stabilizers, multiple slice intersections, and noncompact orbit volume require separate treatment.

The Frobenius frame metric in Eq.~\eqref{eq:pullback_metric} is not invariant under the right action $U_i\mapsto U_i g$ for general $g\in\mathrm{GL}^+(K_q)$, so a quotient metric cannot be inferred from that formula. It must be declared on the quotient or obtained from a suitable invariant ambient metric. Background covariance is a further condition on a redundant-field expansion~\cite{DeWitt1967,Abbott1981,Abbott1982}. Noncompact redundant directions require a quotient or a valid gauge-fixing construction; active noncompact directions require an integrable combined action and reference measure. Haar measure by itself supplies no such normalization.

A two-particle-irreducible extension would introduce a symmetric bilocal source and a second Legendre transform, producing an action $\Gamma_{\mathrm{2PI}}[\varphi,G]$ whose covariance stationarity condition is $\delta\Gamma_{\mathrm{2PI}}/\delta G=0$~\cite{CornwallJackiwTomboulis1974}. A gauge-consistent realization for the present configuration model remains an unresolved construction. It would become relevant if fluctuations change the saddle or if soft modes invalidate the nondegenerate Laplace approximation.

Exact elimination can generate interactions outside the original ansatz, including higher-order region terms, nonlocal kernels, and measure contributions; symmetry constrains the allowed terms but does not require every allowed coefficient to be nonzero~\cite{Wilson1974}. Closure is therefore a conditional approximation to be assessed through the oscillation-to-KL control of Eq.~\eqref{eq:rg_closure} and the separate force discrepancy of Eq.~\eqref{eq:rg_force_error}. Correlated state inference, local configuration fluctuations, and a self-consistent two-point approximation address different errors and do not automatically form nested families or a monotone sequence of evidence bounds. A concrete comparison must specify each normalized model or approximation family and report its evidence gap, correlations, and closure residuals. No spacetime, gauge-field kinetic term, or physical mass follows merely from adopting one of these fluctuation constructions.
